\section{Localization by quadratic curvature}
\label{sec:localization}

We work with one of the active patterns in \eqref{eq:src-11}, on a
normalized segment with root $R$.  Its original column lists, including their
private coordinates, are fixed throughout this section.  Write
\[
 P=\lceil\eta^{-\alpha}\rceil,
\]
where the positive absolute constant $\alpha$ will be chosen in
Section~\ref{sec:completion}.  All estimates below are uniform in the number
of summation depths.  In particular, $\log(2P)=O(k)$, whereas a factor
\[
 G_0=\exp\bigl(\log^{O(1)}(2k)\bigr)
\]
is $P^{o(1)}$.  We distinguish such factors from fixed powers of $P$.
Further attempts can begin at normalized subroots, keeping the original
basic arguments and inheriting their column lists.

The aim is the localization into current groups in
Proposition~\ref{prop:localization}, up to its stated summable error.
A graph first records scale-dependent compatibility of quadratic
curvatures.  Decreasing colors, decreasing graph distances, and
increasing vertex lists will give nested cumulative masks.  The group
masks also retain earlier central uses: the later square estimates need
separation between distinct groups of the same color on overlapping
whole cells that can still lead to an actual use.  These are the usable
endpoints defined below.  The final transfer identity prepares the
two-side estimates of Section~\ref{sec:sides}; the linear component of
those side states is constructed in Section~\ref{sec:linearization}.

\subsection{Tags, an ambient graph, and short paths}

Put $c=(-1,3,-3,1)$.  A vertex is an occurrence of an original basic
column, in any slot.  Its \emph{normalized curvature tag}
$p_v=(\theta_v,C_v,\lambda_v)$ is defined by writing its quadratic chart
phase in slot $i$ as
\[
 c_i\bigl(\lambda_v y^2+(\theta_v y)^T C_v(\theta_v y-n)\bigr).
\]
Thus the curvature data of an original column in slot $i$ are divided by
$c_i$.  A linear modulation is not part of its tag.  The vertex lists
$\mathcal V(I)$ increase on downward paths and have cardinality at most
$D=D_{\rm base}$, enlarging the original base bound if necessary.  They are
ordered by admission, with fixed tie breaking.  Repeated tags can have
different vertices and private labels.

Every additional main or helper column constructed below is attached to
one of these original vertices.  In slot $j$ it has precisely that vertex's
normalized tag, with phase multiplier $c_j$, a bounded smooth chart
amplitude, and an unrestricted linear modulation.  The single-column
complexities remain at most $\exp(k^{O(1)})$, independently of the number
of insertions.  In particular, these insertions do not add vertices to the
graph.  The later linearized modes belong to a separate catalog and do not
add graph vertices either.

At depth $s$, let $r$ be the interval length.  Use the ambient graph of
\emph{all} hypothetical curvature tags satisfying the fixed basic bounds,
not just the currently present vertices.  Identify equal tags at distance
zero.  Join two tags $v,w$ when their difference data
\[
 V=r\binom{\theta_v}{\theta_w},\qquad
 C=\operatorname{diag}(C_v,-C_w),\qquad
 \Lambda=r^2(\lambda_v-\lambda_w)
\]
admit the splitting in \eqref{eq:src-2}, with heights and size bounds at
most $H_{\rm ed}$.  The graph distance is denoted $d_s(v,w)$ and can be
infinite.  The threshold $H_{\rm ed}=\exp(k^{O(1)})$ will be fixed after
an a priori count of the main and helper catalogs.  Refining one depth
halves $V$ and quarters $\Lambda_*$, without changing the rational fast
space.  Consequently edges persist and $d_s$ is nonincreasing in $s$.
These distances depend only on the tags and absolute depth, not on the
location, birth cell, or availability of the vertices.

\begin{lemma}[Separation of different tags]
\label{lem:tag-separation}
Fix an upper bound $\exp(k^{O(1)})$ for the complexities of the single
columns under consideration.  Given
$\tau_2>0$ with $\log(1/\tau_2)\le k^{O(1)}$, one can choose
$H_{\rm ed}=\exp(k^{O(1)})$ so that the following statements hold.
If $d_s(v,w)>1$, the physical inner product of two same-slot pure columns
of these tags, normalized on any interval of depth $s$, has absolute value
at most $\tau_2$.  If the two columns are in distinct slots, their
four-linear form has norm at most $\tau_2$ against two arbitrary
normalized $L^2(I)$ arguments in the other slots.
\end{lemma}

\begin{proof}
For the inner product, smooth the interval indicator with $L^1$ error
much smaller than $\tau_2$ and derivative bounds polynomial in
$\tau_2^{-1}$.  The physical columns have bounded supremum norm.  A
larger inner product would therefore give a Fourier coefficient of their
conjugate product above a fixed multiple of $\tau_2$ after recentering and
rescaling.  Lemma~\ref{lem:weyl} gives \eqref{eq:src-2} for the difference
data.  The fixed factor $c_i$ in a same-slot product only changes the
allowable height by a fixed power.  Choosing $H_{\rm ed}$ above this
bound proves the first assertion.

For the second assertion, call the source slots $i,j$ and the spectator
slots $a,b$.  In centered, rescaled spectator coordinates $(z,w)$, the
linear parts of the source coordinates are
$r(\xi_i z+\zeta_i w)$ and $r(\xi_j z+\zeta_j w)$, where
\[
 \xi_l=\frac{b-l}{b-a},\qquad \zeta_l=\frac{l-a}{b-a}.
\]
All four coefficients for $l=i,j$ are nonzero fixed rationals.  The
identities $\sum_l c_l l^u=0$ for $u=0,1,2$ give
\[
 c_i\xi_i\zeta_i=-c_j\xi_j\zeta_j=:\kappa\ne0.
\]
First clip each spectator at a sufficiently large inverse polynomial in
$\tau_2$.  The discarded part has $L^1$ norm at most the inverse clipping
height, since its $L^2$ norm is bounded.  The absolute form estimate with
exponents $1,2,\infty,\infty$ bounds this error.  We can therefore suppose
that the spectators have both bounded $L^2$ norm and supremum norm at
most an inverse polynomial in $\tau_2$.

Use the return calculation from the proof of Lemma~\ref{lem:weyl} for
the displayed difference data.  Choose an integer period clearing the
chart and interpolation denominators.  For a small return
$Vh=m+d_h$, with $m$ divisible by this period, translate $z$ by $h$ and
translate the source lift indices by $\xi_i m_v,\xi_j m_w$.
The amplitude and kernel errors are bounded by
$\abs{h}+\abs{d_h}$ times a quasipolynomial in the available complexity.
The resulting modulation in $w$ is
\[
 \kappa\Omega(h,m),\qquad
 \Omega(h,m)=(2\Lambda+V^T C V)h-V^T(C-C^T)m.
\]
Indeed the mixed quadratic contribution in each source contains
$c_l\xi_l\zeta_l$; factoring out $\kappa$ leaves exactly the difference
of the two curvature data.  The modulation in $z$ can be included in
the first spectator and has no effect on its $L^2$ norm.

Suppose the original correlation is larger than $\tau_2$.  Applying the calculation to $\ell(h,m)$, for a sufficiently
large quasipolynomial number of integers $\ell$, produces correlations
of size at least a fixed multiple of $\tau_2$ between one fixed bounded
kernel $K_1(z,w)$ and tests $f_\ell(z)g_\ell(w)$.  Here the
$f_\ell$ have bounded $L^2$ norms, and
\[
 g_\ell(w)=g(w)e\bigl(\ell\kappa\Omega(h,m)w\bigr)
\]
up to a common bounded multiplier.  The error remains small by choosing
the return threshold sufficiently small before choosing these multiples.
If $\abs{\Omega(h,m)}$ exceeds a fixed constant, the separated
exponentials, multiplied by the clipped spectator, are Bessel on the
fixed box, with bound controlled by that spectator's squared supremum
norm.  Cauchy--Schwarz in $z$ gives
\[
 \sum_\ell\left|\int K_1(z,w)f_\ell(z)g_\ell(w)\,dz\,dw\right|^2
 \lesssim
 \int\sum_\ell\left|\int K_1(z,w)g_\ell(w)\,dw\right|^2dz.
\]
The right side is bounded by the same Bessel constant times
$\norm{K_1}_2^2$, while the left side grows with the number of returns.
This is a contradiction.  We have proved a uniform bound for
$\Omega(h,m)$ on every return below a quasipolynomially small threshold.

The final, algebraic part of the proof of Lemma~\ref{lem:weyl} now
applies to these difference data: the alternating return identity gives
an isotropic lattice space, successive minima give a rational fast
projection and a bounded slow component, and one return in each
coordinate direction bounds the effective quadratic form.  This part
uses only the just-established return bound.  It yields
\eqref{eq:src-2} at quasipolynomial height, contradicting $d_s(v,w)>1$
when $H_{\rm ed}$ is chosen sufficiently large.
\end{proof}

\begin{lemma}[Endpoint bounds along a short path]
\label{lem:curvature-path}
Suppose two tags can be joined, at one scale, by a path of length
$L\le\operatorname{poly}(P,k)$.  Their difference data satisfy
\eqref{eq:src-2} in their own endpoint coordinates, with
$\log H'\le\operatorname{poly}(P,k)$.
\end{lemma}

\begin{proof}
List the path tags as $v_0,\ldots,v_L$, and put all their horizontal
coordinates in one direct sum.  For each edge choose its fast space
$U_e$ from the definition of an edge.  Let $U$ be the subspace of the
joint coordinate space whose restriction to each consecutive pair
belongs to $U_e$.  These are rational linear equations.  Clearing
denominators and applying determinant bounds to one maximal independent
system shows that the orthogonal projection onto $U$ has logarithmic
height polynomial in $L$, the dimensions, and $\log H_{\rm ed}$.
The same determinant estimate bounds the norm of a right inverse on the
range of this system.

Let $V$ be the concatenation of the normalized horizontal vectors.
Each defining equation has bounded error on $V$, because the edge slow
component is bounded.  The right-inverse bound therefore shows that
$V^F$, the orthogonal projection of $V$ onto $U$, satisfies
$\norm{V-V^F}\le H_1$, where
$\log H_1\le\operatorname{poly}(P,k)$.  On an edge, both its chosen
fast component and the restriction of $V^F$ lie in $U_e$, and their
difference $\Delta$ is bounded by this new budget.  Since $U_e$ is
isotropic for the edge skew form, replacing its fast component by this
restriction changes the effective correction by
\[
 \operatorname{sym}(\Delta^T C_{{\rm sk},e}V_e)
 =\operatorname{sym}(\Delta^T C_{{\rm sk},e}V_{s,e}),
\]
which is bounded at the new budget.  In particular, no unbounded fast
coordinate enters this error.

Sum the edge effective quadratic forms using these common joint fast
and slow components.  The diagonal symmetric terms, the skew mixed
terms, and the original $\lambda$ terms telescope, leaving only the two
endpoints.  Let $U'$ be the image of $U$ under endpoint restriction.
It is rational with the same type of height bound.  For two vectors of
$U$, summing their edge skew pairings cancels all intermediate
coordinates; each edge pairing is zero.  Hence $U'$ is isotropic for
the endpoint difference skew form.  The restricted joint fast vector
lies in $U'$ and is within the new bounded error of the endpoint $V$.
Replacing it by the orthogonal projection of endpoint $V$ onto $U'$
changes it by a bounded vector of $U'$.  The preceding isotropy
calculation again bounds the change in effective curvature.  These are
all the assertions of \eqref{eq:src-2} in endpoint coordinates.
\end{proof}

\subsection{Colors with controlled changes}

The following construction is on original vertices only.  The colors
are ordered integers in a fixed range, and each vertex's color is
nonincreasing along a downward path.  A group will consist of vertices
of one current color together with additional fixed winning data.

Set
\[
 d_0=\lceil1000\log(2D)\rceil,\qquad
 \mathcal L=d_0^{30},\qquad H_{\max}=100P(\mathcal L+1),
\]
and choose the two distance scales
\[
 b_2=100(H_{\max}+1),\qquad
 b_1=d_0^{10}b_2(H_{\max}+1).
\]
These definitions fix the available color range before the coloring
is performed.

In each of $d_0$ independent primary rounds, give every original vertex
$z$ a geometric delay $D_z$, truncated to $[0,d_0]$, with untruncated
tail $\Prob(D_z\ge n)=2^{-n}$.  Arrays for different births and rounds
are independent.  At a present vertex $v$, the top score is
\[
 g=\max_{z\in\mathcal V(I)}
       \{D_z-\lfloor d_s(z,v)/b_1\rfloor\}.
\]
A round is \emph{marked} when it has a unique winner $z$ whose score
exceeds every other score by at least $5$.  Use the first marked primary
round.

Let $i$ be the admission index of its winner and let $j\ge i$ be the
current list size.  If $j=i$, use an equality cohort and the prefix
through $i$.  Otherwise there is a unique smallest dyadic block of
positive indices whose lower half contains $i$ and whose upper half
contains $j$.  Throughout this upper-half cohort freeze $W$ to be the
prefix through that lower half.  The winner $z$ belongs to $W$.
A cohort is specified by this block, and its length level has only
$O(\log(2D))$ possibilities on a path.

For this primary round and frozen context define, at every absolute
depth $s'$ and for every hypothetical tag $a$,
\[
 f_a(s')=\max_{w\in W\setminus\{z\}}
             (b_1D_w-d_{s'}(w,a)).
\]
An empty maximum is $-\infty$.  These profiles are defined before and
after the births of their arguments; availability is not used in this
definition.  They are nondecreasing in depth.  Put
\[
 T_l=(g-3)b_1+\frac{l b_1}{2d_0},\qquad
 U=(g-2)b_1,\qquad 1\le l\le d_0,
\]
and write $C_l(x)=\min(U,\max(T_l,x))$, including the value at
$-\infty$.  The frozen pseudometrics are
\[
 d^{(l)}(a,b)=\sup_{s'\in\Z}
          |C_l(f_a(s'))-C_l(f_b(s'))|.
\]
The supremum over \emph{all absolute depths} is part of the definition.
Thus these metrics are static within the frozen context, and
$d^{(l+1)}\le d^{(l)}$.

Let $R_l=10l(d_0+1)b_2$.  At $v$ choose the least $l<d_0$ for which the
balls with respective metrics and radii
\[
 B_{d^{(l)}}(v,R_l),\qquad
 B_{d^{(l+1)}}(v,R_{l+1}),
\]
counted on currently present vertices, have the same floor logarithm to
base $3/2$.  Call this common integer $n$.  Such an $l$ exists: the balls
are nested as $l$ increases, are nonempty, and their cardinalities are
at most $D$, while $d_0$ exceeds the possible range of their floor
logarithms.

Finally use $d_0$ independent high rounds of geometric delays, independent
of the primary arrays.  The high arrays can be shared among all frozen
contexts.  In the high round use the score defined by
$d^{(l+1)}/b_2$ in place of $d_s/b_1$ and the same gap condition.
Use the first marked high round, with winner $a$ and score $g_2$.

The color metadata are the primary round and score $g$, the equality
case or separating length level, $l,n$, and the high round and score
$g_2$.  At failed markings use fixed placeholders.  Initialize each vertex's change counter to $d_0^8$ at admission.
Whenever any of these data, or either winning identity, changes,
decrement this counter.  Order counter and metadata
lexicographically, with the counter first, and embed the resulting
finite ordered set into $\{1,\ldots,\mathcal L\}$.  A successful
\emph{group} additionally records $z,a,W$ and the exact cohort block.
It does not record the changing list size within an upper half or the
identity of the marked central vertex.

\begin{lemma}[Color budget and exceptional roots]
\label{lem:color-budget}
For sufficiently large starting $k$, the preceding construction uses at
most $\mathcal L$ colors, with at most $d_0^8$ changes per vertex on a
path.  Delays can be chosen so that the first cells where a required
primary or high marking fails have arbitrarily small fixed total length
relative to $R$.  They can be sent to fresh attempts.  All statements
below are on the remaining paths, where every needed marking exists.
\end{lemma}

\begin{proof}
At a present vertex a round's top score lies in $\{0,\ldots,d_0\}$:
the vertex itself is a competitor of distance zero.  Individual primary
competitor scores only increase with depth.  At a fixed top score a
marked winner can occur in at most one episode, with a fixed identity.
Once a rival has destroyed its gap, that rival's score cannot decrease;
a different marked winner at the same top value would have the old
winner as a rival of that value.  Thus the first marked primary round,
its score, and its winner have at most $5d_0^2$ episodes.

Within a primary episode there are at most $d_0$ cohorts.  Within one
cohort all profile metrics are fixed.  The vector of floor log-sizes of all the
balls has at most $2d_0^2$ epochs, including its initial value, since
its coordinates increase and each has at most $O(\log D)$ possible
log-sizes.  Within each such epoch the selected $l,n$ are fixed.  In that epoch the high metric is fixed,
so high scores increase only by admission of competitors.  The same
marked-episode argument gives at most $5d_0^2$ episodes of selected high
data.  Altogether there are at most $50d_0^7\le d_0^8$ changes when
$k$ is large.  Seven metadata coordinates, each with at most $d_0+2$
states, and the change counter require fewer than
$256d_0^{15}<d_0^{30}$ ordered states.  Decrementing the counter at each
change makes the assigned color nonincreasing even if the other
metadata increase.

For completeness, a fixed untruncated score situation is marked with
probability at least $2^{-5}$.  Choose a canonical winner in every delay
outcome, resolving ties deterministically, and increase only its delay
by $5$.  The resulting outcome has a unique winner with gap at least
$5$.  This map is injective, because its image identifies the increased
coordinate; geometric probabilities decrease by exactly $2^{-5}$.
Truncation changes this probability by at most $D2^{-d_0}$.  In
particular it remains at least $1/64$ for large $k$.

There are only polynomially many situations to check on a fixed
spatial path.  For the primary stage cap
$\lfloor d_s(z,v)/b_1\rfloor$ at $d_0+5$; larger values cannot affect
any top or gap decision.  The capped pair distances change only
$O(D^2(d_0+1))$ times, by monotonicity.  Counting the vertices to test
gives $O(D^3(d_0+1))$ fixed primary situations.  Conditional on all
primary arrays, the primary, cohort, and ball epochs number at most
$O(Dd_0^5)$ across vertices.  Each has at most $D$ present prefixes,
so at most $O(D^2d_0^5)$ high situations require checking.  These
situations are fixed before revealing the independent high arrays.
The probability that all $d_0$ rounds fail at any such situation is
at most $(1-1/64)^{d_0}$.  The union bound tends to zero with $D$,
since $d_0\ge1000\log(2D)$.

Integrate this pathwise failure probability over $R$.  Some realization
of the arrays has a small measure of failed paths.  The first failure
cells are disjoint, and their total length is at most that measure.
This proves the claimed exceptional-root bound without a bound on the
number of vertices across different spatial branches.
\end{proof}

\subsection{Historical ball masks and usable endpoints}

Fix a successful group $G$.  For an integer pad $1\le h\le H_{\max}$,
include a present vertex $b$ if, for some historical central vertex
$v_0$ assigned this same color and group at or above the current cell
on its ancestor path,
\begin{equation*}
 \begin{split}
 b_1D_z-d_s(z,b)&\ge(g-1)b_1-h,\\
 d^{(l+1)}(b,v_0)&\le h,\qquad
 d^{(l)}(b,v_0)\le10R_lh.
 \end{split}
 \tag{12}\label{eq:src-12}
\end{equation*}
Denote this vertex mask by $\mathcal M_{G,h}$.  It is empty before its
first central assignment.  Within the cohort it increases in depth:
the score increases, the metrics are frozen, and historical witnesses
are retained.  After the cohort, where no further actual central use of
this group occurs on that branch, retain the eligible vertex set by
inheritance whenever a nested operator is needed.  Births on other
branches supply no historical data to the current branch.

An \emph{actual use} of a group requires a current central vertex
assigned to it.  A \emph{usable endpoint} for it is any cell with a
nonempty pad mask and some actual use of that group at or below the cell.  Thus a
usable endpoint need not itself be central.  Nonempty usable endpoints
are still in their cohorts: an inherited mask after a cohort has no
later actual use there.

\begin{lemma}[Mask inclusion, closure, and separation]
\label{lem:historical-masks}
The following assertions hold for every allowed pad.
\begin{enumerate}
\item At an actual use with current central vertex $v$, every present
vertex $b$ with $d_s(b,v)\le h$ lies in $\mathcal M_{G,h}$.
\item At an actual use, every present graph neighbor of a vertex in
$\mathcal M_{G,h}$ belongs to $\mathcal M_{G,h+1}$ if $h<H_{\max}$.
Every vertex admitted to a mask has, already at admission, a graph path
to $z$ of length at most $\operatorname{poly}(P,k)$.
\item For distinct groups of the same color, vertices belonging to their
pad masks on overlapping usable endpoints are at graph distance greater
than $1$ at the finer endpoint.
\end{enumerate}
\end{lemma}

\begin{proof}
At a marked primary center $v$ at depth $s$, each rival's true score is
at most $(g-5)b_1$, since its quantized score is at most $g-5$.
Consequently
\[
 f_v(s')\le(g-5)b_1<T_l\quad(s'\le s).
\]
The winning true score is greater than $(g-1)b_1$.  If $d_s(b,v)\le h$,
the winner score at $b$ loses at most $h$, and
$f_b(s)\le(g-5)b_1+h<T_l$.  Both profiles are below the clipping range
at every earlier depth.  At all later depths the graph distance remains
at most $h$, so their true profiles, and then their clipped profiles,
differ by at most $h$.  Thus both profile distances from $b$ to $v$ are
at most $h$, which proves the first assertion with witness $v_0=v$.

Next compare a historical central witness $v_0$ with the current center
$v$ of the same group.  Their common high winner $a$ has delay at most
$d_0$, and the marked top score $g_2$ is nonnegative.  Hence
\[
 d^{(l+1)}(a,v_0),\ d^{(l+1)}(a,v)<(d_0+1)b_2.
\]
The earlier low ball at $v_0$, on the vertices present when it was
chosen, is contained in the current high ball at $v$: use
$d^{(l+1)}\le d^{(l)}$ and
$R_l+2(d_0+1)b_2<R_{l+1}$.  This earlier low ball and the current low
ball each have at least $(3/2)^n$ vertices, whereas the current high
ball has fewer than $(3/2)^{n+1}$ vertices.  The two low balls intersect.
The triangle inequality in the low metric gives
\[
 d^{(l)}(v_0,v)\le2R_l.
\]
For a pad-$h$ vertex $b$ it follows that
$d^{(l)}(b,v)\le A_h:=10R_lh+2R_l$.  Our constants give
\[
 A_h<T_{l+1}-T_l-10\qquad(h\le H_{\max}).
\]
For example the gap on the right before subtracting $10$ is
$\tfrac12d_0^9b_2(H_{\max}+1)$, whereas
$A_{H_{\max}}\le10d_0(d_0+1)b_2(10H_{\max}+2)$.
Since the clipped profile of $v$ at the use is $T_l$, this proves the
\emph{true} rival bound
\[
 f_b(s)\le T_l+A_h<T_{l+1}-10.
\]

Let $b'$ be a present graph neighbor of $b$ at this use.  Its winner score loses
at most one.  Both $f_b$ and $f_{b'}$ are below $T_{l+1}$ at all depths
up to $s$, and their difference is at most one at all depths after $s$.
Thus $d^{(l+1)}(b,b')\le1$.  For the low metric, after $s$ the distance
from $b'$ to $v_0$ costs at most $10R_lh+1$.  Before $s$, its clipped
profile is between $T_l$ and $T_l+A_h+1$, while the clipped profile of
$v_0$ is between $T_l$ and $T_l+2R_l$.  Their difference is at most
$A_h+1\le10R_l(h+1)$.  All three predicates of
\eqref{eq:src-12} therefore hold at pad $h+1$.  Finally, at admission
the first predicate gives
\[
 d_s(z,b)\le(d_0+1)b_1+H_{\max},
\]
which is polynomial in $P,k$ and proves the path bound.

It remains to prove the usable-endpoint assertion.  Let $I,J$ be
overlapping endpoints for distinct same-color groups $G,G'$, with $J$
the finer endpoint, and choose an actual-use certificate $u$ for
$G'$ at or below $J$.  Let $b,b'$ be the respective mask vertices.  Suppose,
for a contradiction, that $d_{\operatorname{depth}(J)}(b,b')\le1$.
This distance bound persists to $u$.  The vertex $b'$ remains in its mask at $u$.
No future use of $G$ on this same branch is required.

First suppose their primary winners $z,z'$ differ.  Their colors give
the same primary round, score $g$, and equality case or separating
length level.  At a fixed length level, upper-half index cohorts on an
increasing path either use the same dyadic block or increasing disjoint
blocks.  In either case the earlier winner belongs to the later frozen
prefix $W'$.  This is immediate also for equality cohorts.  The earlier
winner score at $b$ persists, so adjacency at $u$ gives
\[
 f'_{b'}(u)\ge(g-1)b_1-h-1.
\]
This contradicts the true rival bound for $G'$ at $u$, which is less
than $T_{l+1}-10\le(g-5/2)b_1-10$, since $b_1$ is much larger than
$H_{\max}$.

Now suppose $z=z'$.  Its fixed admission index and the common
separating level determine the same cohort block; its prefix $W$ is inherited on the
overlapping path.  The metrics coincide, and distinct groups must have distinct high
winners $a,a'$.  Historical witnesses $v_0,v'_0$ for the two masks lie
on the common ancestor path of $J$ and can be ordered along it.  At the
later historical witness the earlier high winner remains available in
the same high round and metric.  The earlier winner's true score is
greater than $(g_2-1)b_2$ at its own witness, whereas at the later
witness it is at most $(g_2-5)b_2$.  Therefore
\[
 d^{(l+1)}(v_0,v'_0)>4b_2>2H_{\max}+1.
\]
On the other hand, at the certificate $u$ the true rival bound puts
$f_{b'}(u)$ below $T_{l+1}-10$, and adjacency puts $f_b(u)$ below
$T_{l+1}-9$.  Both profiles were below the high clipping range at all
earlier depths, and they differ by at most one at all later depths.
It follows that $d^{(l+1)}(b,b')\le1$.  The two mask predicates then
imply
\[
 d^{(l+1)}(v_0,v'_0)\le h+1+h'\le2H_{\max}+1,
\]
a contradiction.  This proves separation, including the case where the
two endpoints have later certificates on different branches.
\end{proof}

\begin{lemma}[Bessel separation of usable endpoint spaces]
\label{lem:usable-bessel}
Fix a slot and a color.  For each of finitely many groups $G$ let $Z_G$
be a finite sum of available main pure columns, with constant
coefficients on \emph{entire} usable endpoint cells whose pad masks
contain their tags.  Different terms may use arbitrary endpoint times,
and a label may be used repeatedly.  Suppose the number of groups and
available column labels on a spatial path, and the inverse private
weights, are bounded by $\exp(k^{O(1)})$.  Choosing $\tau_2$ sufficiently
small after these bounds gives
\[
 \norm{\sum_G Z_G}^2\lesssim\sum_G\norm{Z_G}^2.
\]
The constant has no dependence on the number of endpoint times or on
the depth of a future use certificate.
\end{lemma}

\begin{proof}
For a fixed label $v$ in $G$, combine all its occurrences into the scalar
coefficient function
\[
 a_{v,G}(y)=\sum_I a_{v,G,I}\1_I(y).
\]
The private coordinate of label $v$ in $Z_G(y)$ is precisely
$\epsilon_v^{1/2}a_{v,G}(y)$.  Thus, if $W_0$ bounds inverse square-root
private weights,
\[
 \sum_v|a_{v,G}(y)|^2\le W_0^2\norm{Z_G(y)}^2.
\]
This uses the combined coefficients, so any cancellation between repeated
uses of a label is retained.

Fix labels $v,w$ in distinct groups.  Their tags have a deterministic
last separating depth: by monotonicity, the depths where
$d_s(v,w)>1$ form an initial interval.  Cap this interval at the last
depth in the finite calculation if needed, and call its endpoint $s_*$.
If two supporting endpoint cells $I,J$ overlap, their intersection is
the finer cell.  Lemma~\ref{lem:historical-masks} says that its depth
is at most $s_*$.  Therefore every nonzero term in
\[
 a_{v,G}\overline{a_{w,G'}}
  =\sum_{I,J}a_{v,G,I}\overline{a_{w,G',J}}\1_{I\cap J}
\]
is measurable on the grid at depth $s_*$.  The whole product is
measurable on that grid.  Endpoint terms finer than $s_*$ have no
overlap with any term of the other label and contribute zero.  If no
separating depth occurs within the root, there are no overlaps at all.

Apply Lemma~\ref{lem:tag-separation} on every grid cell.  The two private
labels cannot coincide on an overlap, since equal tags have distance
zero.  Hence
\[
 |\ip{Z_G}{Z_{G'}}|
 \le\tau_2\sum_{v,w}\int|a_{v,G}(y)a_{w,G'}(y)|\,dy.
\]
Pointwise Cauchy--Schwarz in groups and labels bounds the sum of all
these cross terms by $\tau_2$ times a fixed power of the pathwise count
and private-weight bounds, multiplied by
$\sum_G\norm{Z_G}^2$.  Taking $\tau_2$ smaller than the reciprocal of
this factor proves the assertion.  There is no factor counting
repetitions or endpoint times.
\end{proof}

The whole-cell hypothesis of Lemma~\ref{lem:usable-bessel} will be used
also for adjoint ranges.  If an output is restricted to finer use cells,
the adjoint puts that restriction on the input side of its endpoint
projection.  The projection then again produces constant column
coefficients on the entire endpoint cell.  Thus earlier endpoint cells
are retained in full whenever they have a downstream certificate;
spatial restrictions are not treated as arbitrary restrictions of the
range spaces in the lemma.

\subsection{Same-tag tests and finite projection catalogs}

\begin{lemma}[A uniform same-tag test]
\label{lem:same-tag}
For each original vertex $v$ there is a fixed template family in every
slot, consisting of bounded pure columns of normalized tag $v$ and
unrestricted linear modulations, with the following property.  If a
pure source column of tag $v$ in slot $i$ has form norm $t_0>0$ against
a physical residual in slot $j$ and two normalized $L^2(I)$ spectators,
then the residual has normalized moment at least $t_0/H_0$ against a
template column of the same tag in slot $j$.  Here
$H_0\le\exp(k^{O(1)})$ is independent of all insertion ranks and of the
linear modulations.  The source can be either an original column or
any column in these template families.
\end{lemma}

\begin{proof}
Fix a nonnegative compact smooth partition function $\rho$ with
$\sum_{n\in\Z^d}\rho(u-n)=1$.  Its support and Gevrey bounds can be
chosen within the basic budget.  Choose a sufficiently large integer
period $P_0$, fixed for this vertex, accommodating the supports of both
its original amplitude and these partition functions.  The templates
have amplitude
\[
 \rho(d)e(\ell^T d/P_0),\qquad d=\theta_v y-n,
\]
where $\ell$ ranges over residue vectors modulo $P_0$, in addition to
an arbitrary external linear modulation.  Their physical supremum norms
are bounded by one because the partition weights are nonnegative.

Insert this partition in the other three slots and use lifts
$d_l=\theta_v y_l-n_l$ in all four slots.  Extract the matching phase
$c_l(\lambda_vy_l^2+(\theta_vy_l)^TC_vd_l)$ from each input, putting
its conjugate into the corresponding input multiplier.  On the
progression,
\[
 d_l=(1-l)d_0+ld_1+k_l,
\]
with bounded integer carries $k_l$.  The common quadratic terms cancel
because $\sum_l c_l l^u=0$ for $u=0,1,2$.  The phase left by the
carries is
\[
 \sum_l c_l(\theta_v y_l)^T C_v k_l,
\]
which is a sum of linear modulations.  There are quasipolynomially many
possible carries within the fixed support bounds.

To separate the lifts, impose each specified carry equality by a smooth
function of $d_l-(1-l)d_0-ld_1$ that is one at that integer and zero at
all other possible integers.  Such a function is exact here, since its
argument on the progression is integer valued.  Expand these coupled
smooth functions in Fourier series on a box of period $P_0$ containing
the bounded lifts.  After multiplication by the partition functions,
every nonsource factor has the required template amplitude, possibly
with an unrestricted character index.  Write an index as
$\ell=P_0q+r$; then
\[
 e(\ell^Td/P_0)=e(q^T\theta_v y)e(r^Td/P_0),
\]
since $q^Tn$ is integral.  Thus only the residue $r$ enters the template,
and the other part is an allowed linear modulation.  The source factor,
after cancellation of its phase, periodizes a smooth horizontal
function; expand that function in ordinary torus characters.  Its
remaining dependence on $y_i$ is only a linear modulation.  Derivatives
of sufficiently high order, polynomial in the dimension, bound all these
absolute Fourier coefficient sums by a quasipolynomial in the fixed
basic budget.  For a template source, reducing its own character to a
residue gives the same bound.  There is no iterative increase of
nonlinear complexity.

It remains to prove the assertion for one separated term.  Express
$y_i,y_j$ in the two spectator coordinates $y_a,y_b$.  Separate the
smooth rescaled kernel by an absolutely summable Fourier expansion in
these coordinates; source linear modulations can be absorbed into the
two spectator tests.  The target coordinate has the form
\[
 y_j=\alpha y_a+\beta y_b,
\]
with both $\alpha,\beta$ nonzero fixed rationals.  Zero-extend the
demodulated target, consisting of the physical residual multiplied by
the conjugate of its bounded template, from $I$ to a fixed larger
interval and Fourier-expand it there.  Its coefficients are normalized
moments of the residual with same-tag templates and arbitrary linear
modulations.  If their supremum is $M$, the resulting integral is
bounded by
\[
 M\sum_m|A_mB_m|
 \le M\Bigl(\sum_m|A_m|^2\Bigr)^{1/2}
         \Bigl(\sum_m|B_m|^2\Bigr)^{1/2}.
\]
Here $A_m,B_m$ are the spectator coefficients at frequencies affine in
$m$, with nonzero spacings proportional to $\alpha,\beta$.
These families are Bessel on the fixed integration intervals.
Multiplication of a spectator by its bounded demodulating template
preserves its $L^2$ bound.  The last display is therefore at most a
fixed constant times $M$.  Fourier expansions can first be finite;
$L^2$ convergence and the same Bessel estimate pass to the limit.
Summing the absolutely convergent separated terms gives the form bound
$H_0M$.  A form norm larger than $t_0$ consequently forces a target
moment at least $t_0/H_0$, after enlarging $H_0$ by a fixed factor if
necessary to avoid an unattained supremum.
\end{proof}

Besides the ball masks, define
\[
 \mathcal U_{m,h}(I)=
 \{v\in\mathcal V(I):d_s(v,\{w:\text{color}(w)\le m\})\le h\}.
\]
They increase in depth, because distances decrease, colors do not
increase, and vertices are inherited.  They also have the immediate
neighbor-closure property with pad $h+1$.

Give the original basic labels rank zero.  In a slot, a projection of
rank $r$ and a vertex mask means the ordinary cell-local orthogonal
projection onto all available columns of ranks at most $r$ whose tags
belong to that mask.  These are projections of vector functions with
constant column coefficients on the cell.  Newly inserted columns have
private weight one and a new private coordinate.  A positive smeared
filter averages $P$ such projections at the same rank, with pads
$h+p$, $1\le p\le P$.  At a fixed cell these projections are nested
in $p$; at a fixed pad they are nested in depth.  Their averages and
complements are contractions, although products of them need not be
projections.

For the cumulative filter at a split $m$, use rank
$\mathcal L-m+1$ and base pad $10P(\mathcal L-m+1)$.  At a final group,
use four ball filters of ranks $\mathcal L+j$, $1\le j\le4$, in their
insertion order, with base pads
\[
 20P(\mathcal L+1)+10(j-1)P.
\]
All these pads, including one-neighbor enlargements used below, are
within $H_{\max}$.

We require the following finite closure of the catalogs.  At a target
rank $r$ and one of its individual masks, apply its residual projection
to a required history of length $O(\log\mathcal L)$ on a fixed basic
$B$ or $U$ argument.  Test that physical residual against every pure
available lower-rank source whose tag lies in the target mask and two
arbitrary normalized $L^2$ spectators.  The required form norm is at
most
\begin{equation*}
 \frac{\tau_3}{D\sqrt{n_{r-1}}}
        \norm{\text{base of target history}}_{2,I},
 \qquad \tau_3=D^{-C_\tau}.
 \tag{13}\label{eq:src-13}
\end{equation*}
Here $n_r$ is an a priori bound for all labels through rank $r$ on a
path, and $C_\tau$ is a sufficiently large fixed exponent.  The required
histories include all individual pad choices of the products and
complements used in the localization and transfer below.  Thus the
same tests hold for averaged histories.  Columns, once inserted, are
inherited in descendants.

There is also a helper rank $\mathcal L+5$, constructed after the side
states in Section~\ref{sec:sides} are fixed.  Its sources are main
columns only.  Its additional bases are the whole side states, with
at most $\exp(k^{O(1)})$ bases per active cell over slots, groups, and
patterns.  These bases are fixed independently of helper insertions;
they are not individual frequency expansions.  They may have extra
private coordinates, all chosen distinct from fresh helper coordinates.
The same history length and tolerance \eqref{eq:src-13} are required.

\begin{lemma}[Catalog closure and its budget]
\label{lem:catalog-closure}
The main closure and then the helper closure can be performed at every
active cell with path counts satisfying
$\log n_{\max}\le k^{O(1)}$.  The exponent is fixed independently of
$H_{\rm ed}$ and of subsequent linearization precision.  All fresh
columns have the template complexities of Lemma~\ref{lem:same-tag}.
\end{lemma}

\begin{proof}
Hold the original graph vertices, masks, and named histories fixed.
The numbers of masks and histories at one cell are
$\exp(k^{O(1)})$: there are polynomially many original-vertex contexts,
polynomially many ranks and colors, $P$ pad choices per factor, and
$O(\log\mathcal L)$ factors.  This count uses names of contexts and
masks, not their numeric distance values or the number of inserted
columns.  The helper bases have the stipulated same count.

Upon a failed test, Lemma~\ref{lem:same-tag} supplies a bounded physical
column with moment at least the threshold in \eqref{eq:src-13} divided
by $H_0$.  Insert a fresh copy at the target rank.  Its new private
coordinate is zero in the old tested state and in the old projection.
Thus that physical moment equals the vector inner product with the old
orthogonal residual.  If $P_{\rm old},P_{\rm new}$ are the target
projections before and after insertion and $X$ is the old tested state,
then
\[
 \norm{(P_{\rm new}-P_{\rm old})X}_{2,I}
 \gtrsim \frac{\tau_3}{D\sqrt{n_{r-1}}H_0}
        \norm{\text{base}}_{2,I}.
\]
Indeed the fresh vector column has uniformly bounded norm, and
projection onto its component orthogonal to the old range gives this
inequality.  If the base has zero norm no test fails.

Although $X$ can change when other projections change, its squared
increments can still be charged.  Fix a named history, a target mask,
and a trial sequence of at most $T$ insertions.  Every individual
positive projection in that history is nested in the insertion index.
Expand complements and use the binary-prefix argument of
Lemma~\ref{lem:compression}.  The target increments are orthogonal
rows; the right-hand history is evaluated at the old insertion index.
The sum of their squared actions on this changing history is at most
\[
 (C\log(2T))^{O(\log\mathcal L)}\norm{\text{base}}_{2,I}^2.
\]
This is precisely the row version of that binary argument, and it does
not require treating a changing projected state as a fixed vector.

It follows that rank-$r$ insertions at a cell are bounded by
\[
 n_{r-1}\exp(k^{C_1})
             (C\log(2T))^{O(\log\mathcal L)}.
\]
The exponential factor includes the configuration count, the squared
inverse tolerance, and $H_0^2$.  It does not include the number of
possible triggering source columns.  For a fixed target history, any
triggering source produces a fresh moment in that history's old
residual; the resulting target increment is charged to the same
orthogonal row sequence regardless of that source's identity.

There are at most $M$ active cells on a path.  Absorb this factor, and
fixed factors for ranks and slots, in $\exp(k^{C_1})$.  Starting with
a basic bound $n_0$, choose
\[
 n_r/n_{r-1}=\exp(k^{C_2})
\]
with $C_2$ sufficiently large, and set the global trial cap to the
resulting $n_{\max}$, with spare fixed factors.  Since there are only
$\mathcal L+5=k^{O(1)}$ ranks, $\log n_{\max}=k^{O(1)}$.
With $T=n_{\max}$ the binary-prefix factor is much smaller than the
chosen ratio for large $k$.  Thus no trial cap is attained.  This
argument also accommodates simultaneous changes of histories caused
by insertions at other ranks: all the underlying positive projections
remain nested along the common insertion sequence.  It proves
termination of the main closure, and the identical argument proves
termination of the later helper closure on its fixed bases.

All constants used in this count were determined by the basic template
bounds and the number of named graph contexts.  They do not depend on
$H_{\rm ed}$.  We may therefore first fix these bounds, then choose
$\tau_2$ sufficiently small for Lemma~\ref{lem:usable-bessel} and all
pair-mismatch errors, and finally choose $H_{\rm ed}$ by
Lemma~\ref{lem:tag-separation}.  This order is not circular.
\end{proof}

We may summarize the resulting main count, coefficient, inverse private
weight, and active-multiplicity bounds by $D_{\rm m}$, where
$\log D_{\rm m}\le k^{O(1)}$.  The dimensions of individual tags and
unions of two tags obey the same polynomial bound, independently of
$P$.  Whenever an outer primal range of rank at most $r$ is expanded,
its coefficient $\ell^1$ norm is at most
\[
 D\sqrt{n_r}\,\norm{\text{its input}}_{2,I}.
\]
This follows from the private-coordinate lower bound and contraction
of the projection.  It is the reason for the normalization in
\eqref{eq:src-13}; the full column coefficient sum will not be charged
in the main estimates.

\subsection{Binary localization}

\begin{proposition}[Localization into current groups]
\label{prop:localization}
For an active pattern in \eqref{eq:src-11}, there are
$\operatorname{poly}(k)$ fixed search leaves such that, up to total
absolute length-weighted error at most $Ce^{-k}|R|$, each leaf contributes
\begin{equation*}
 \sum_{G\ \mathrm{current}}
 H_I(Z_{0,G,s},Z_{1,G,s},Z_{2,G,s},Z_{3,G,s}),\qquad
 Z_{l,G,s}=\overline L_{l,G,s}\mathcal B_{l,s}z_l.
 \tag{14}\label{eq:src-14}
\end{equation*}
Here $\overline L_{l,G,s}$ is one positive smeared ball filter, and
$\mathcal B_{l,s}$ is a product of cumulative filters or complements
of length $O(\log\mathcal L)$, independent of $G$.  The product and
slot choices are fixed by the search leaf, uniformly over cells.  A
group is summed only at an actual current central use.  The minimum
with $\eta$ in \eqref{eq:src-11} may be applied separately per leaf.
\end{proposition}

\begin{proof}
Maintain a range of colors $[\ell,h]$, an upper primal guard slot, and
when $\ell>1$ a lower residual guard slot.  Initially the range is
$[1,\mathcal L]$; use any of the at least three basic coarse arguments
as the upper guard and take no lower guard.  At a midpoint
$m=\lfloor(\ell+h)/2\rfloor$, choose a third slot distinct from the
guards and split it by the smeared cumulative filter at $m$ and its
complement.  In the positive term that slot becomes the upper guard,
the former upper guard is released, and the range becomes $[\ell,m]$.
In the residual term it becomes the lower guard, the former lower
guard is released, and the range becomes $[m+1,h]$.  A guard is not
postcomposed before it is released.  Use a deterministic slot-choice
rule for each branch of this binary search, even if some colors are
empty at a particular cell.  There are $O(\log\mathcal L)$ splits per
path and at most $\mathcal L$ leaves.

Consider a leaf with color $k'$.  Its upper guard is either the initial
basic coarse argument or has outer cumulative rank
$\mathcal L-k'+1$ and the corresponding pad.  Expand this one outer
primal support, keeping its average when present.  For the initial basic
guard, set the upper pad to zero.  Call a source tag
\emph{old} if it lies within this upper pad of a vertex of color at
most $k'-1$.  If $k'>1$, the lower residual guard has rank
$\mathcal L-k'+2$ and a base pad larger by at least $10P$.  Every old source
therefore lies in all its residual masks and has lower rank.  Equation
\eqref{eq:src-13} kills it, with arbitrary normalized $L^2$ spectators.
For $k'=1$ there are no old tags.

Every remaining upper source is near a current vertex of color $k'$.
In the initial basic case this is the source's own original vertex;
in the other case it follows from membership in the cumulative mask
and exclusion of old tags.  Let $G$ be the current group of such a
central vertex.  The upper pad is smaller than the first ball pad, so
Lemma~\ref{lem:historical-masks} puts the source in every needed mask
of that ball filter.  The source cannot match two distinct groups:
it would then violate same-color mask separation at the current cell.
Moreover it is separated from the primal sources of every other
current group's ball filter.

Choose a third slot, distinct from the upper and lower guards, and
insert its first-rank ball projection, summed over current groups.
For the matching group, the difference from the identity is bounded
by \eqref{eq:src-13}, using the upper source.  For any other group,
expand the new primal ball source and apply
Lemma~\ref{lem:tag-separation}.  Old-source terms on either side are
still bounded using the untouched lower residual guard.  This yields
an identity, with small error, in which the third slot has pure support
in one group's first ball mask.

Keep that third slot as the source, and postcompose the other three
slots with their own larger-rank, larger-pad ball filters in order.
The first ball source has lower rank than each new target, and its
mask is contained in all the new target masks.  Thus
\eqref{eq:src-13} bounds each residual insertion.  The third slot is
left unchanged.  The result is \eqref{eq:src-14}; each slot's inner
cumulative history depends only on the search path, and only its outer
ball filter depends on the group.

We specify why the numerical errors are summable without spending a
full coefficient bound in a main term.  Each test using a rank-$r$
residual expands only an outer lower-rank source.  Its coefficient
bound is at most $D\sqrt{n_{r-1}}$ times the local input norm, canceled
by the denominator in \eqref{eq:src-13}.  Basic local sizes and all
same-time contractions are bounded.  The remaining counts of active
cells on a path, leaves, current groups, and filter-insertion steps are fixed
powers of $D$; choosing $C_\tau$ sufficiently large makes their total
length-weighted contribution at most $e^{-k}|R|$.  Pair-mismatch
errors can additionally contain fixed powers of the full main count,
but $\tau_2$ was chosen after that count and makes the same bound hold.
No expansion of a full column coefficient sum is retained in
\eqref{eq:src-14}.  Finally,
$\min(\eta,\sum_j a_j)\le\sum_j\min(\eta,a_j)$ for $a_j\ge0$,
so the minimum can be distributed over the fixed leaves as claimed.
\end{proof}

For a basic coarse input $z_l=B_l$, the histories before its ball filter
are fixed signed sums of products of averaged nested column projections
on $F_l$, shared across groups.  Their product lengths are
$O(\log\mathcal L)$, so Lemma~\ref{lem:compression} gives losses at most
$G_0$.  Before an attempt root, ball outputs can be set to zero.
Between actual uses the positive masks and histories are still evaluated
on their inherited endpoint cells whenever needed for nesting.

\subsection{Transferring complementary filters to a side}

\begin{lemma}[Outgoing filters and near-band defects]
\label{lem:transfer}
At a use of \eqref{eq:src-14}, suppose slots $a,b$ are fixed side states
$S_a,S_b$, chosen independently of the helper catalog.  Suppose the
other two slots have whole arguments $Z_{c,G,s},Z_{d,G,s}$ with basic
inputs $B_c,B_d$.  Expand every factor $1-\overline Q$ in slots $c,d$
as the difference of the identity and a positive average.  In each
resulting term, all positive filters on these two slots can be removed,
one at a time, by transferring corresponding helper averages to slot $a$.
The main term at the end is
\[
 H_I(\mathcal D_{a,s}S_a,S_b,B_c,B_d),
\]
where $\mathcal D_{a,s}$ is a product of helper averages of length
$O(\log\mathcal L)$.

For a step transferring a filter from $c$ to $a$, let $Q_{c,p}$ be
its individual outgoing projections, $1\le p\le P$, and let $Q_{a,p}$
be the helper-rank projections with the same masks and pads.
For $l=a,c$, put
\[
 Q_{l,0}=0,\qquad Q_{l,P+1}=1,\qquad
 A_l^p=Q_{l,p}-Q_{l,p-1},\qquad
 \overline Q_l=P^{-1}\sum_{p=1}^P Q_{l,p}.
\]
Let $x_c$ be the inner history on $B_c$ after deleting the outgoing
filter.  In the other slots, let $x_a=\mathcal D_{a,s}S_a$ use the helper
averages already transferred, let $x_b=S_b$, and let $x_d$ be the
whole remaining history on $B_d$.  The exact difference
\[
 H_I(x_a,x_b,\overline Q_c x_c,x_d)
 -H_I(\overline Q_a x_a,x_b,x_c,x_d)
\]
is
\begin{equation*}
 \sum_{p,q=1}^{P+1}\frac{p-q}{P}
 H_I(A_a^p x_a,x_b,A_c^q x_c,x_d).
 \tag{15}\label{eq:src-15}
\end{equation*}
Only $|p-q|\le2$ need be retained.  The discarded far terms at each
cell and group have total absolute size at most
$D^{-C_3}\norm{S_a}_{2,I}\norm{S_b}_{2,I}$, for any prescribed fixed
$C_3$, after the indicated choices of numerical accuracies.
\end{lemma}

\begin{proof}
Expand each complementary factor $1-\overline Q$ as a difference of
identity and a positive average.  The number of resulting products is
at most $2^{O(\log\mathcal L)}=\operatorname{poly}(k)$.  In any such
product, remove the larger rank among the current outermost filters of
$c,d$, regarding an unfiltered basic input as rank zero.  Positive
ranks are distinct: each cumulative rank was inserted in exactly one
slot at one search split, and the four final ball ranks are distinct.
Thus the other complementary source always has strictly lower rank
than the filter being removed.

Suppose the next filter is on $c$, and use the projections and arguments
defined in the statement.  The $A_l^p$ are orthogonal band projections
at the current cell and sum to the identity.

On band $p$, the positive average $P^{-1}\sum_{u=1}^P Q_{l,u}$ has
weight $(P-p+1)/P$, including weight zero at $p=P+1$.  Therefore the
difference in the statement is exactly Equation~\eqref{eq:src-15},
with the arguments written in role order.
This proves the algebraic identity, including the terminal bands.

If $p>q+2$, the lower $c$ band is a difference of primals, so expand
its pure sources.  Their rank is a main rank below the helper rank,
and their tags belong to both helper masks needed to write the higher
$a$ band as a difference of residuals.  The helper version of
\eqref{eq:src-13}, applied to the history on the fixed base $S_a$,
bounds the resulting terms.  The two other arguments are arbitrary
$L^2$ spectators.

If $q>p+2$, instead expand the outer primal source of $x_d$, using the
basic expansion if it is already unfiltered.  This is a main source
of rank strictly below the outgoing rank on $c$.  The lower $a$ band
is primal.  Let $\mathcal M_p$ be the outgoing vertex mask at its pad
$h+p$.  Classify a pure source of $x_d$ as near if its tag is within
one edge of some vertex in $\mathcal M_p$, and as far otherwise.
For a near source, neighbor closure gives membership in
$\mathcal M_{p+1}$, hence in both masks needed for the higher $c$
residual band.  Apply the main test with the entire helper band
$A_a^p x_a$ as one $L^2$ spectator.  For a far source, expand the
helper primal: every pair of source tags is separated, so
Lemma~\ref{lem:tag-separation} applies.  Neighbor closure is the
triangle inequality for a cumulative mask and
Lemma~\ref{lem:historical-masks} at the current actual use for a
ball mask.  The gap
$q-p>2$ leaves at least this one-pad margin.  The main version of
\eqref{eq:src-13}, already imposed on the inner history $x_c$, then
bounds it.  Thus helper insertions need not alter any main catalog.

In the first direction the coefficient bound of the main source is
canceled by the helper test's denominator.  In the second direction
the lower-rank main source coefficient is canceled by the main test's
denominator; the helper primal serves only as an $L^2$ spectator.
For pair mismatch, expanding the two primal supports can cost fixed
powers of the helper and main counts.  Those counts were anticipated
before $\tau_2$ was selected.  There are only $O(P^2)$ pairs of bands
and polynomially many transfer positions.  Since $\log P=O(k)$, these
factors are absorbed by the prior accuracy choices.  All filters are
contractions, so the remaining norm factors are
$\norm{S_a}_{2,I}\norm{S_b}_{2,I}$ and bounded complementary basic
sizes.  This proves the stated far-error estimate.  Averages in
$\mathcal D$ are bounded directly by their tested averaged histories
and do not cost the number of pad choices.

Retain the near bands and forward only the transferred main term to
the next step.  Iterating deletes all complementary filters and gives
the asserted final main term.  The number of helper factors remains
$O(\log\mathcal L)$.  Restoring $B_c,B_d$ to $Y_c,Y_d$ in that main
term uses the full pair estimate; every resulting difference contains
an additional basic side $U$.  These are the three-side terms estimated
in Section~\ref{sec:sides}.
\end{proof}

We record the range convention for the retained defects.  Keep their
band sums together at each fixed value of $q-p$.  If the outgoing filter
is a ball filter, its final complementary band
$A_c^{P+1}=1-Q_{c,P}$ is split into the group-independent identity and
the positive ball term.  All other complementary band endpoints have
usable ball ranges.  If the outgoing filter is cumulative, its outer
ball has already been removed, and the whole complementary band stack
is shared across groups.  Likewise slot $d$ either has one positive
ball filter outermost or is shared.  These conventions concern the
estimation of defects, not the transfer of its main term, and will
permit the joint square bounds of Section~\ref{sec:sides}.
